\documentclass[11pt]{article}
\usepackage[margin=1in]{geometry}
\usepackage{amsmath,amssymb,amsthm,bm}
\usepackage{array,booktabs,longtable}

\newtheorem{theorem}{Theorem}
\newtheorem{lemma}{Lemma}
\newtheorem{definition}{Definition}
\newtheorem{assumption}{Assumption}

\title{\textbf{SCS --- Internal Technical Spec}\\[2pt]
\large Score-only sparse-electrode attack on EEG recognition}
\author{}
\date{Internal. Edit this first; the manuscript follows it.}

\begin{document}
\maketitle
\vspace{-3em}

%====================================================================
\section{Question}

\paragraph{Question.} Under \emph{score-only} access (the adversary sees class
logits, not gradients or weights), how few \emph{complete electrode traces} must be
controlled to flip an EEG subject-recognition decision using only EEG-band waveforms,
and \emph{what quantity governs that number}? The question is open because dense,
white-box, and channel-time attacks do not measure whole-channel cost, and because the
\emph{mechanism} behind a small count --- not just its value --- is unknown.

%====================================================================
\section{Setup and notation}

All objects are typed before use. The trial index is carried on every per-trial
object; it is dropped only for \emph{operators} whose value does not depend on which
trial they act on (the margin and the waveform map).

\paragraph{Notation typography rule.}
Scalars use ordinary italic or Greek letters, e.g.\ $c,k,t,\rho,\Delta,\nu$ and the
declared scalar scores $F_c,\widehat F_c$. Vectors use bold lowercase, e.g.\
$\mathbf{u}_j$ and $\mathbf{s}(\mathbf{E})$. Matrices or multichannel trials use bold
uppercase, e.g.\ $\mathbf{X}_i$, $\mathbf{E}_i$, and $\mathbf{A}_i^{(k)}$. Sets use
calligraphic uppercase, e.g.\ $\mathcal{D}$, $\mathcal{C}$, $\mathcal{B}$, and
$\mathcal{S}^{(k)}$. In particular, $s_c(\mathbf{E})$ is the scalar channel indicator,
while $\mathcal{S}(\mathbf{E})$ is the corresponding channel set; $e_{i,c}^{(k)}(t)$ is
a scalar waveform value, while $\mathbf{E}_{\mathcal{S}^{(k)},\mathbf{A}_i^{(k)}}$ is the
full perturbation matrix.

\paragraph{Indexing and dependency invariant.}
Before any equation is accepted, every symbol must declare what it depends on, its
size/shape, and the set or source it comes from. Indexing is one case of this rule: if the
right-hand side depends on trial $i$, channel $c$, atom $m$, support size $k$, probe draw
$j$, probe sample $q$, or optimization step $\ell$, then the left-hand side must expose
that dependency. Table~\ref{tab:dependency-ledger} is the gate used for SCS notation.

\paragraph{Data and model.}
The dataset is $\mathcal{D}=\{(\mathbf{X}_i,y_i)\}_{i=1}^N$ with trials
$\mathbf{X}_i\in\mathbb{R}^{C\times T}$, where $C$ is the number of electrode channels
and $T$ is the number of samples. Let $\mathcal{C}:=\{1,\dots,C\}$ be the channel index
set. Labels are $y_i\in[K]:=\{1,\dots,K\}$, and $i\in[N]$. The recognizer is a score map
$f:\mathbb{R}^{C\times T}\to\mathbb{R}^{K}$ with prediction
$\hat{y}(\mathbf{X})=\arg\max_{r\in[K]} f_r(\mathbf{X})$, using a fixed deterministic
tie-break. The clean predicted identity of trial $i$ is $y_0^i=\hat{y}(\mathbf{X}_i)$.

\paragraph{Data-use accounting.}
The data are not one undifferentiated pool. The model-training split
$\mathcal{D}_{\mathrm{train}}$ is used only to fit the recognizer $f$. The validation
split $\mathcal{D}_{\mathrm{val}}$ is used to report clean accuracy/EER and to form the
clean-correct pool
$\mathcal{D}_{\mathrm{cc}}=\{(\mathbf{X},y)\in\mathcal{D}_{\mathrm{val}}:
\hat y(\mathbf{X})=y\}$. Stage~1 uses a separate probe set
$\mathcal{D}_{\mathrm{probe}}=\{\mathbf{Z}_q\}_{q=1}^{n}$ to estimate channel leverage.
Stage~2 evaluates attacks on
$\mathcal{D}_{\mathrm{attack}}\subseteq\mathcal{D}_{\mathrm{cc}}$. The final experiment
must report both $n=|\mathcal{D}_{\mathrm{probe}}|$ and
$|\mathcal{D}_{\mathrm{attack}}|$; ASR and $k^\star$ are computed only over
$\mathcal{D}_{\mathrm{attack}}$, not over all validation windows.

\begin{center}
\footnotesize
\begin{tabular}{@{}>{\raggedright\arraybackslash}p{0.22\textwidth}
>{\raggedright\arraybackslash}p{0.28\textwidth}
>{\raggedright\arraybackslash}p{0.42\textwidth}@{}}
\toprule
Object & Current HBN 122-subject artifact & Use \\
\midrule
$\mathcal{D}_{\mathrm{train}}$ & 785 passive-task recordings; 55,665 windows &
train EEGConformer subject recognizer \\
$\mathcal{D}_{\mathrm{val}}$ & 487 active-task recordings; 44,694 windows &
report clean validation accuracy and biometric EER \\
$\mathcal{D}_{\mathrm{cc}}$ & 35,042 clean-correct validation windows &
candidate pool from which attack trials are drawn \\
$\mathcal{D}_{\mathrm{attack}}$ & 122 clean-correct windows, capped by
\texttt{attack\_max\_samples=122} & report ASR, attacked accuracy, and
first-success channel count $k^\star$ \\
$\mathcal{D}_{\mathrm{probe}}$ & not yet regenerated for the population Stage~1 artifact &
estimate $F_c$ and $\widehat F_c$ for the shared channel ranking \\
\bottomrule
\end{tabular}
\end{center}

\paragraph{Margin (operator).}
For any $\mathbf{X}\in\mathbb{R}^{C\times T}$, $y\in[K]$,
\begin{equation}
M(\mathbf{X},y)=f_y(\mathbf{X})-\max_{r\neq y} f_r(\mathbf{X}).
\end{equation}
The clean decision on trial $i$ is correct iff $M(\mathbf{X}_i,y_0^i)>0$; a non-targeted
attack on trial $i$ succeeds once $M(\mathbf{X}_i+\mathbf{E}_i,y_0^i)<0$.

\paragraph{Perturbation, support, budget.}
The perturbation chosen for trial $i$ is $\mathbf{E}_i\in\mathbb{R}^{C\times T}$ (a
per-trial solution, hence indexed). The channel-support indicator is
\[
s_c(\mathbf{E})=\mathbf{1}\{\|\mathbf{E}_{c,:}\|_2>0\},\qquad c\in\mathcal{C},
\]
with vector $\mathbf{s}(\mathbf{E})\in\{0,1\}^{C}$. The activated channel set, when needed,
is $\mathcal{S}(\mathbf{E})=\{c\in\mathcal{C}:s_c(\mathbf{E})=1\}$, and the
channel-count cost is $\sum_{c\in\mathcal{C}}s_c(\mathbf{E})$. The support is not fixed
in advance --- $\mathcal{S}(\mathbf{E}_i)$ is whatever channels the solution activates.
Each perturbation
obeys the peak-ratio cap
$\|\mathbf{E}_i\|_\infty\le r_{\max}\,\|\mathbf{X}_i\|_\infty$, $r_{\max}\in(0,1)$.

\paragraph{Waveform basis.}
Before either stage runs, we build a fixed waveform dictionary
$\mathcal{B}=\{b_m\}_{m=1}^{r}$ for one EEG window. It is not learned from the model and
not resampled during the attack. It is constructed from known signal-processing choices:
the window length $T$, sampling rate $f_s$, Hanning taper $h(t)$, an EEG-band frequency
grid, a grid of bump centers, and low-order trend degrees. Each atom
$b_m:\{0,\dots,T-1\}\to\mathbb{R}$ is normalized to unit norm,
$b_m=\tilde b_m/\|\tilde b_m\|_2$. The atom indices are partitioned into three families,
$\mathcal{M}_{\mathrm{osc}}\cup\mathcal{M}_{\mathrm{hat}}\cup
\mathcal{M}_{\mathrm{trend}}=[r]$, with sizes
$r=r_{\mathrm{osc}}+r_{\mathrm{hat}}+r_{\mathrm{trend}}$. For oscillatory atoms, use a
fixed product grid of frequencies and phases:
$\mathcal{F}=\{f_a^{\mathrm{osc}}\}_{a=1}^{R_f}\subseteq[f_{\min},f_{\max}]$ and
$\mathcal{G}_{\varphi}=\{\varphi_b\}_{b=1}^{R_\varphi}\subset S^1$. Each
$m\in\mathcal{M}_{\mathrm{osc}}$ corresponds to one pair
$(\alpha(m),\beta(m))\in[R_f]\times[R_\varphi]$. These are fixed dictionary coordinates,
not random variables sampled during the attack:
\begin{align}
\text{(osc)}\quad   &\tilde b_m(t)=h(t)\cos(2\pi f_{\alpha(m)}^{\mathrm{osc}} t/f_s+\varphi_{\beta(m)}),
&&m\in\mathcal{M}_{\mathrm{osc}};\\
\text{(hat)}\quad   &\tilde b_m(t)=h(t)\max\!\big(1-|t/(T-1)-\mu_m|/w,\,0\big),
&&m\in\mathcal{M}_{\mathrm{hat}};\\
\text{(trend)}\quad &\tilde b_m(t)=h(t)\big(s(t)^{d_m}-\overline{s^{d_m}}\big),\ s(t)=\tfrac{2t}{T-1}-1,
&&m\in\mathcal{M}_{\mathrm{trend}}.
\end{align}
For an active channel set $\mathcal{S}\subseteq\mathcal{C}$ and coefficient matrix
$\mathbf{A}\in\mathbb{R}^{C\times r}$ with one row per global EEG channel, the induced
perturbation operator uses the set indicator directly:
\begin{equation}\label{eq:param}
(\mathbf{E}_{\mathcal{S},\mathbf{A}})_{c,t}=
\mathbf{1}\{c\in \mathcal{S}\}\sum_{m=1}^{r} a_{c,m}\,b_m(t).
\end{equation}
Here $\mathbf{A}$ is an operator argument. We keep it as a $C\times r$ matrix instead of a
local $|\mathcal{S}|\times r$ matrix so row $c$ always means channel $c$. In Stage~2, the
optimized coefficient matrix is trial-and-budget dependent, written $\mathbf{A}_i^{(k)}$,
so the coefficient on trial $i$, channel $c$, atom $m$, and budget $k$ is
$a_{i,c,m}^{(k)}$.

\paragraph{Sub-Gaussian.}
A real random variable $Y$ with mean $\mu$ is \emph{$\nu$-sub-Gaussian} if
$\mathbb{E}\,e^{\lambda(Y-\mu)}\le e^{\lambda^2\nu^2/2}$ for all $\lambda\in\mathbb{R}$;
then $\Pr(|Y-\mu|\ge\varepsilon)\le 2e^{-\varepsilon^2/(2\nu^2)}$, and the mean of $n$
i.i.d.\ such variables is $(\nu/\sqrt n)$-sub-Gaussian.

%====================================================================
\section{Threat model objective}

For each trial the ideal attack is the per-trial combinatorial program
\begin{equation}\label{eq:ideal}
\min_{\mathbf{E}_i\in\mathbb{R}^{C\times T}} \sum_{c\in\mathcal{C}}s_c(\mathbf{E}_i)
\quad\text{s.t.}\quad M(\mathbf{X}_i+\mathbf{E}_i,y_0^i)<0,\;
\|\mathbf{E}_i\|_\infty\le r_{\max}\|\mathbf{X}_i\|_\infty .
\end{equation}
The minimizer $\mathbf{E}_i^\star$ and its activated channel set
$\mathcal{S}(\mathbf{E}_i^\star)$ depend on $\mathbf{X}_i$.
SCS is a black-box heuristic for \eqref{eq:ideal}, not a certificate of its optimum.

%====================================================================
\section{Method: two-stage SCS}

\subsection{Stage 1: population channel ranking}
Let $P_{\mathrm{probe}}$ denote the distribution of probe trials available to the
score-only attacker. Fix a finite probe set
$\mathcal{D}_{\mathrm{probe}}=\{\mathbf{Z}_q\}_{q=1}^{n}$ with
$\mathbf{Z}_q\stackrel{\mathrm{i.i.d.}}{\sim}P_{\mathrm{probe}}$. The set
$\mathcal{D}_{\mathrm{probe}}$ is the sample used to estimate the ranking; it is not the
distribution itself. Also fix a bank
$\mathcal{P}=\{\mathbf{u}_1,\dots,\mathbf{u}_{J}\}$ of $J$ sign vectors
$\mathbf{u}_j\in\{-1,+1\}^{r}$. The bank is sampled once, e.g.\ with independent
Rademacher entries, and then held fixed for all channels and probe trials.

For channel $c$ and probe direction $\mathbf{u}_j$, the unscaled single-channel probe is
\[
(\widetilde{\mathbf{E}}_{c,j})_{c',t}=
\begin{cases}
\rho\sum_{m=1}^{r}u_{j,m}b_m(t), & c'=c,\\
0, & c'\neq c,
\end{cases}
\]
where $\rho>0$ controls the probe amplitude. For any probe-distribution trial
$\mathbf{X}$, define the capped probe
$\mathbf{E}^{\mathrm{probe}}_{\mathbf{X},c,j}=\mathrm{Cap}_{\mathbf{X}}
(\widetilde{\mathbf{E}}_{c,j})$, a deterministic rescaling to satisfy the peak-ratio cap.
For the finite sample $\mathbf{Z}_q$, this becomes
$\mathbf{E}^{\mathrm{probe}}_{\mathbf{Z}_q,c,j}$. Thus the basis $\{b_m\}$, channel $c$,
and amplitude $\rho$ are fixed; the waveform randomness comes from the sampled sign
vector $\mathbf{u}_j$, and the cap can make the final amplitude depend on the probe trial
being evaluated.
\begin{definition}[leverage]\label{def:lev}
The leverage of channel $c$ on any probe-distribution trial $\mathbf{X}$ is the average
single-channel margin drop over the probe bank,
\[
g_c(\mathbf{X})=M(\mathbf{X},\hat y(\mathbf{X}))-\frac{1}{J}\sum_{j=1}^{J}
M\!\big(\mathbf{X}+\mathbf{E}^{\mathrm{probe}}_{\mathbf{X},c,j},\,\hat y(\mathbf{X})\big).
\]
It is an \emph{average} (a sample statistic), not a best case, so its population mean is
estimable by a sample mean.
\end{definition}
The population and empirical leverages and the induced order are
\[
F_c=\mathbb{E}_{\mathbf{X}\sim P_{\mathrm{probe}}}[g_c(\mathbf{X})],\qquad
\widehat F_c=\frac1n\sum_{q=1}^{n} g_c(\mathbf{Z}_q),\qquad
\widehat\pi=\text{argsort}_{\downarrow}\,\widehat F,\qquad
\pi^\star=\text{argsort}_{\downarrow}\,F,
\]
where $F_c$ is an expectation over the population distribution $P_{\mathrm{probe}}$ and
$\widehat F_c$ is the sample mean over the finite set $\mathcal{D}_{\mathrm{probe}}$.
All sorting uses a fixed deterministic tie-break. Here $\widehat\pi$ is the empirical
order used by SCS and $\pi^\star$ is the population order used only for analysis. For a
target support size $k<C$, the \emph{frontier gap} is
$\Delta=F_{\pi^\star_k}-F_{\pi^\star_{k+1}}\ge 0$, and $\nu$ is the sub-Gaussian proxy of
$g_c(\mathbf{X})$ for $\mathbf{X}\sim P_{\mathrm{probe}}$
(Assumption~\ref{as:subg}).

\paragraph{Assumptions for Claim~1.}
The ranking bound rests on three assumptions. They are stated after the probe objects,
leverage function, and expectation are defined, so every symbol in the assumptions has a
declared dependency.
\begin{assumption}[unbounded scores]\label{as:logits}
The oracle returns logits in $\mathbb{R}^K$, not probabilities. Hence concentration
must use a sub-Gaussian tail, not a bounded-range (Hoeffding) tail.
\end{assumption}
\begin{assumption}[independent probes]\label{as:indep}
The finite probe trials $\mathbf{Z}_1,\ldots,\mathbf{Z}_n$ in
$\mathcal{D}_{\mathrm{probe}}$ are independent draws from $P_{\mathrm{probe}}$; in
practice this is approximated by taking one window per distinct recording.
\end{assumption}
\begin{assumption}[sub-Gaussian leverage]\label{as:subg}
For each channel $c\in\mathcal{C}$, the random variable $g_c(\mathbf{X})$ with
$\mathbf{X}\sim P_{\mathrm{probe}}$ is $\nu$-sub-Gaussian about its mean $F_c$. If it is
only finite-variance, a Bernstein bound gives the same $1/\Delta^2$ rate with a larger
constant.
\end{assumption}

\subsection{Stage 2: per-trial attack in fixed order}
Fix a maximum channel budget $B<C$. For each $k\in\{1,\ldots,B\}$, the active channel set
is the universal empirical prefix
$\mathcal{S}^{(k)}=\{\widehat\pi_1,\dots,\widehat\pi_k\}$
($|\mathcal{S}^{(k)}|=k$), \emph{identical across trials}: only the coefficients
$\mathbf{A}_i^{(k)}$ and the realized count $k_i^\star$ are per-trial. For trial $i$ at
budget $k$, the perturbation on selected channel $c$ is
\[
e_{i,c}^{(k)}(t)=\sum_{m=1}^{r} a_{i,c,m}^{(k)}b_m(t),\qquad c\in\mathcal{S}^{(k)}.
\]
The atoms $b_m$ are fixed across trials, but the coefficients $a_{i,c,m}^{(k)}$ are
optimized for the particular trial $\mathbf{X}_i$ and selected set size $k$. Define the
trial-and-budget objective
\[
J_i^{(k)}(\mathbf{A})=M(\mathbf{X}_i+\mathbf{E}_{\mathcal{S}^{(k)},\mathbf{A}},y_0^i).
\]
Stage~2 replaces the random probe coefficients from Stage~1 with optimized
trial-specific coefficients. SPSA refines them using two score queries per iteration.
At iteration $\ell$, with current coefficients $\mathbf{A}_{i,\ell}^{(k)}$, random sign
matrix $\mathbf{U}_{i,\ell}^{(k)}$ supported only on rows in $\mathcal{S}^{(k)}$, and radius
$d_\ell$, the SPSA estimate is
\[
\widehat{\mathbf{G}}_{i,\ell}^{(k)}
=\frac{
J_i^{(k)}(\mathbf{A}_{i,\ell}^{(k)}+d_\ell\mathbf{U}_{i,\ell}^{(k)})
-J_i^{(k)}(\mathbf{A}_{i,\ell}^{(k)}-d_\ell\mathbf{U}_{i,\ell}^{(k)})
}{2d_\ell}\,\mathbf{U}_{i,\ell}^{(k)}.
\]
The update is
\[
\mathbf{A}_{i,\ell+1}^{(k)}
=\mathrm{Proj}_{\mathbf{X}_i,\mathcal{S}^{(k)}}\!\left(
\mathbf{A}_{i,\ell}^{(k)}-q_\ell\widehat{\mathbf{G}}_{i,\ell}^{(k)}
\right),
\]
with Spall schedules $d_\ell=d_0/(\ell+1)^{0.101}$, $q_\ell=q_0/\sqrt{\ell+1}$, and
$\mathrm{Proj}_{\mathbf{X}_i,\mathcal{S}^{(k)}}$ enforcing the peak cap relative to
$\mathbf{X}_i$ and zero rows outside $\mathcal{S}^{(k)}$. After the query budget for support size
$k$ is spent, let $\mathbf{A}_i^{(k)}$ be the best coefficient matrix found for trial
$i$ at that support size.

\paragraph{Metric.} The adversarial trial returned for $\mathbf{X}_i$ at budget $k$ is
\[
\mathbf{X}_{i,\mathrm{adv}}^{(k)}
=\mathbf{X}_i+\mathbf{E}_{\mathcal{S}^{(k)},\mathbf{A}_i^{(k)}}.
\]
The first-success count of trial $i$ is
\[
k_i^\star=\min\{k\in[B]:M(\mathbf{X}_{i,\mathrm{adv}}^{(k)},y_0^i)<0\},
\]
undefined if no $k\le B$ succeeds. The successful set is
$\mathcal{I}^\star=\{i\in[N]:k_i^\star\ \text{is defined}\}$. Over this set we report ASR,
mean/median of
$\{k_i^\star\}_{i\in\mathcal{I}^\star}$, and the prefix-ASR knees $K90/K95/K100$.

\paragraph{Dependency table for SCS.}
\begingroup
\footnotesize
\setlength{\tabcolsep}{3pt}
\renewcommand{\arraystretch}{1.15}
\begin{longtable}{@{}>{\raggedright\arraybackslash}p{0.18\textwidth}
>{\raggedright\arraybackslash}p{0.30\textwidth}
>{\raggedright\arraybackslash}p{0.23\textwidth}
>{\raggedright\arraybackslash}p{0.19\textwidth}@{}}
\caption{Index, dependency, shape, and source ledger for SCS notation.}
\label{tab:dependency-ledger}\\
\toprule
Symbol & Depends on & Size/shape & Source or set \\
\midrule
\endfirsthead
\toprule
Symbol & Depends on & Size/shape & Source or set \\
\midrule
\endhead
$i$ & none & scalar index & $[N]$ \\
$c,c'$ & none & scalar channel index & $\mathcal{C}=\{1,\ldots,C\}$ \\
$t$ & none & scalar time index & $\{0,\ldots,T-1\}$ \\
$m$ & none & scalar atom index & $[r]$ \\
$q$ & none & scalar probe-sample index & $[n]$ \\
$j$ & none & scalar probe-direction index & $[J]$ \\
$k$ & none & scalar support budget & $[B]$, with $B<C$ \\
$\ell$ & none & scalar optimizer step & SPSA iteration set \\
$\mathbf{X}_i$ & $i$ & $C\times T$ real matrix & $\mathcal{D}$ \\
$y_i$ & $i$ & scalar label & $[K]$ \\
$y_0^i$ & $i,\mathbf{X}_i,f$ & scalar predicted label & $[K]$ \\
$\mathcal{D}_{\mathrm{train}}$ & HBN passive-task split & set of labeled trials/windows & subset of $\mathcal{D}$ \\
$\mathcal{D}_{\mathrm{val}}$ & HBN active-task split & set of labeled trials/windows & subset of $\mathcal{D}$ \\
$\mathcal{D}_{\mathrm{cc}}$ & $\mathcal{D}_{\mathrm{val}},f$ & set of clean-correct trials/windows & subset of $\mathcal{D}_{\mathrm{val}}$ \\
$\mathcal{D}_{\mathrm{attack}}$ & $\mathcal{D}_{\mathrm{cc}},\text{attack sample cap}$ & finite attack-evaluation set & subset of $\mathcal{D}_{\mathrm{cc}}$ \\
$f$ & trained model & map $\mathbb{R}^{C\times T}\to\mathbb{R}^K$ & score oracle \\
$M(\mathbf{X},y)$ & $\mathbf{X},y,f$ & scalar margin & $\mathbb{R}$ \\
$\mathbf{E}_i$ & $i,\mathbf{X}_i,y_0^i,f,r_{\max}$ & $C\times T$ real matrix & feasible perturbation set \\
$s_c(\mathbf{E})$ & $\mathbf{E},c$ & scalar indicator & $\{0,1\}$ \\
$\mathbf{s}(\mathbf{E})$ & $\mathbf{E}$ & $C$-vector & $\{0,1\}^{C}$ \\
$\mathcal{S}(\mathbf{E})$ & $\mathbf{E}$ & channel subset & $2^{\mathcal{C}}$ \\
$\mathcal{F}$ & $f_{\min},f_{\max},R_f$ & set of $R_f$ frequencies & $[f_{\min},f_{\max}]$ \\
$\mathcal{G}_{\varphi}$ & $R_\varphi$ & set of $R_\varphi$ phases & $S^1$ \\
$\alpha(m),\beta(m)$ & $m,\mathcal{M}_{\mathrm{osc}},\mathcal{F},\mathcal{G}_{\varphi}$ & pair of scalar grid indices & $[R_f]\times[R_\varphi]$ \\
$b_m$ & $m,T,f_s,\mathcal{F},\mathcal{G}_{\varphi},\alpha,\beta,\text{dictionary grid}$ & length-$T$ unit-norm waveform & $\mathcal{B}$ \\
$\mathbf{A}$ & operator argument & $C\times r$ coefficient matrix & one row per channel in $\mathcal{C}$ \\
$\mathbf{E}_{\mathcal{S},\mathbf{A}}$ & $\mathcal{S},\mathbf{A},\mathcal{B}$ & $C\times T$ perturbation matrix & waveform operator image \\
$P_{\mathrm{probe}}$ & attacker's probe source & distribution over trials & distributions on $\mathbb{R}^{C\times T}$ \\
$\mathbf{Z}_q$ & $q,P_{\mathrm{probe}}$ & $C\times T$ real matrix & $\mathcal{D}_{\mathrm{probe}}$ \\
$\mathcal{D}_{\mathrm{probe}}$ & $n,P_{\mathrm{probe}}$ & set of $n$ trials & finite probe sample \\
$\mathbf{u}_j$ & $j,\text{probe-bank seed}$ & length-$r$ sign vector & $\{-1,+1\}^{r}$ \\
$\mathcal{P}$ & $J,\text{probe-bank seed}$ & set of $J$ sign vectors & probe bank \\
$\mathbf{E}^{\mathrm{probe}}_{\mathbf{X},c,j}$ & $\mathbf{X},c,j,\rho,r_{\max},\mathcal{B}$ & $C\times T$ perturbation matrix & capped single-channel probe \\
$g_c(\mathbf{X})$ & $\mathbf{X},c,\mathcal{P},\rho,r_{\max},\mathcal{B},f$ & scalar margin drop & $\mathbb{R}$ \\
$F_c$ & $c,P_{\mathrm{probe}},\mathcal{P},\rho,r_{\max},\mathcal{B},f$ & scalar population leverage & $\mathbb{R}$ \\
$\widehat F_c$ & $c,\mathcal{D}_{\mathrm{probe}},\mathcal{P},\rho,r_{\max},\mathcal{B},f$ & scalar empirical leverage & $\mathbb{R}$ \\
$\pi^\star$ & $F_1,\ldots,F_C$ & channel permutation & $\mathcal{C}^{C}$ \\
$\widehat\pi$ & $\widehat F_1,\ldots,\widehat F_C$ & channel permutation & $\mathcal{C}^{C}$ \\
$\mathcal{S}^{(k)}$ & $\widehat\pi,k$ & $k$-element channel subset & subsets of $\mathcal{C}$ \\
$\mathbf{A}_{i,\ell}^{(k)}$ & $i,k,\ell,\mathbf{X}_i,y_0^i,\mathcal{S}^{(k)},\text{SPSA seed}$ & $C\times r$ coefficient matrix & SPSA state, zero outside $\mathcal{S}^{(k)}$ \\
$\mathbf{A}_i^{(k)}$ & $i,k,\mathbf{X}_i,y_0^i,\mathcal{S}^{(k)},f,\text{query budget}$ & $C\times r$ coefficient matrix & best Stage~2 state \\
$a_{i,c,m}^{(k)}$ & $i,c,m,k,\mathbf{A}_i^{(k)}$ & scalar coefficient & $\mathbb{R}$ \\
$e_{i,c}^{(k)}(t)$ & $i,c,k,t,\mathbf{A}_i^{(k)},\mathcal{B}$ & scalar waveform value & $\mathbb{R}$ \\
$J_i^{(k)}(\mathbf{A})$ & $\mathbf{A},i,k,\mathbf{X}_i,y_0^i,\mathcal{S}^{(k)},\mathcal{B},f$ & scalar objective value & $\mathbb{R}$ \\
$\mathbf{U}_{i,\ell}^{(k)}$ & $i,k,\ell,\mathcal{S}^{(k)},\text{SPSA seed}$ & $C\times r$ sign/zero matrix & SPSA perturbation direction \\
$\widehat{\mathbf{G}}_{i,\ell}^{(k)}$ & $i,k,\ell,J_i^{(k)},\mathbf{A}_{i,\ell}^{(k)},\mathbf{U}_{i,\ell}^{(k)},d_\ell$ & $C\times r$ matrix & SPSA gradient-like estimate \\
$\mathbf{X}_{i,\mathrm{adv}}^{(k)}$ & $i,k,\mathbf{X}_i,\mathcal{S}^{(k)},\mathbf{A}_i^{(k)},\mathcal{B}$ & $C\times T$ real matrix & attacked trial \\
$k_i^\star$ & $i,B,y_0^i,\{\mathbf{X}_{i,\mathrm{adv}}^{(k)}\}_{k=1}^{B},f$ & scalar or undefined & $\{1,\ldots,B\}$ when successful \\
$\mathcal{I}^\star$ & $\{k_i^\star\}_{i=1}^{N}$ & trial-index subset & subsets of $[N]$ \\
$\Delta$ & $F,\pi^\star,k$ & nonnegative scalar & $\mathbb{R}_{\ge 0}$ \\
$\nu$ & $\{g_c(\mathbf{X}):c\in\mathcal{C},\mathbf{X}\sim P_{\mathrm{probe}}\}$ & nonnegative scalar proxy & sub-Gaussian parameter \\
\bottomrule
\end{longtable}
\endgroup
For example, $e_{i,c}^{(k)}(t)=\sum_m a_{i,c,m}^{(k)}b_m(t)$ passes because the left side
shows the trial, channel, budget, and time dependencies. Writing $e_c(t)$ would fail
because it hides the trial and budget dependence.

%====================================================================
\section{Claim 1 --- channel-ranking stability}

\begin{theorem}\label{thm:rank}
Condition on the fixed dictionary $\mathcal{B}$ and probe bank $\mathcal{P}$. Under
Assumptions~\ref{as:indep}--\ref{as:subg}, fix $\delta\in(0,1)$ and a target size $k<C$
with frontier gap $\Delta>0$. If
\begin{equation}\label{eq:nstar}
n\;\ge\;\frac{8\,\nu^2\,\log(2C/\delta)}{\Delta^2},
\end{equation}
then with probability at least $1-\delta$ the empirical top-$k$ set
$\{\widehat\pi_1,\dots,\widehat\pi_k\}$ equals the population top-$k$ set
$\{\pi^\star_1,\dots,\pi^\star_k\}$.
\end{theorem}

\begin{lemma}[deterministic separation]\label{lem:det}
If $\max_{c\in\mathcal{C}}|\widehat F_c-F_c|<\Delta/2$, the empirical and population top-$k$ sets
coincide.
\end{lemma}
\begin{proof}
Let $a$ be any kept channel (population rank $\le k$) and $b$ any dropped channel (rank
$>k$). By definition of $\Delta$, $F_a\ge F_{\pi^\star_k}\ge
F_{\pi^\star_{k+1}}+\Delta\ge F_b+\Delta$. Then
$\widehat F_a> F_a-\Delta/2\ge F_b+\Delta/2>\widehat F_b$. So every kept channel outranks
every dropped channel empirically; the top-$k$ set is unchanged.
\end{proof}

\begin{proof}[Proof of Theorem~\ref{thm:rank}]
Each $\widehat F_c=\frac1n\sum_{q=1}^{n} g_c(\mathbf{Z}_q)$ is a mean of $n$ independent
(Assumption~\ref{as:indep}) $\nu$-sub-Gaussian (Assumption~\ref{as:subg}) terms, hence
$\Pr(|\widehat F_c-F_c|\ge\varepsilon)\le 2e^{-n\varepsilon^2/(2\nu^2)}$. A union bound
over the $C$ channels gives
$\Pr(\max_{c\in\mathcal{C}}|\widehat F_c-F_c|\ge\varepsilon)\le
2C\,e^{-n\varepsilon^2/(2\nu^2)}$.
Take $\varepsilon=\Delta/2$; the right side is $\le\delta$ iff
$n\ge 8\nu^2\log(2C/\delta)/\Delta^2$, which is \eqref{eq:nstar}. On that event
Lemma~\ref{lem:det} applies, giving top-$k$ set recovery with probability $\ge1-\delta$.
\end{proof}

\paragraph{Insight (the transferable content).} The \emph{form} of \eqref{eq:nstar}, not
the number $n$, is the result: probe cost grows as $\log C$ in the number of electrodes
but as $1/\Delta^2$ in the frontier gap. Closely-ranked electrodes, not many electrodes,
cost queries. A small $k_i^\star$ is not implied by the ranking theorem alone; it is
evidence that the early high-$F_c$ electrodes are sufficient to flip many decisions once
the order is stable.
We therefore treat \eqref{eq:nstar} as a prediction: estimate empirical proxies for
$\nu$ and $\Delta$ from probe data, then check that ranking agreement saturates near the
predicted $n^\star$.

%====================================================================
\section*{Status (internal)}
\small
HBN sweep table/figures present but numbers are from the OLD per-trial greedy; regenerate
under the population attack. Claim~1 confirmation (sweep $n$; ranking agreement vs $n$;
overlay $n^\star$ from empirical proxies for $\nu,\Delta$; error bars) not yet run. BNCI reference
comparison done. Open: whether to headline the law in the abstract; build the
\texttt{notation.tex} \verb|\newcommand| ledger mirroring \S2.

\end{document}
